\documentclass[aspectratio=169,10pt,spanish]{beamer}

\input{../../../_preambulo_beamer.tex}
\input{../../../_comandos_beamer.tex}

\selectlanguage{spanish}

\title{MA1001B - Modelación Estadística para la Toma de Decisiones}
\subtitle{Lección 10: Diagramas de árbol y descomposición del espacio muestral}
\author{}
\institute{Tecnol\'ogico de Monterrey}
\date{}

\begin{document}

\begin{frame}[plain]
  \vspace{1.2cm}
  {\Large\bfseries MA1001B - Modelación Estadística para la Toma de Decisiones\par}
  \vspace{0.7cm}
  {\large Lección 10: Diagramas de árbol y descomposición del espacio muestral\par}
  \vspace{0.7cm}
  {\color{TecRojo}\rule{\textwidth}{1.2pt}}
  \vspace{0.8cm}

  Facultad MA1001B\\[0.3cm]
  Periodo\\[0.3cm]
  Tecnol\'ogico de Monterrey
\end{frame}

\begin{frame}{La pregunta de probabilidad secuencial}
  \begin{alertblock}{Pregunta guía}
    ¿Cómo puede un experimento de dos etapas descomponerse en ramas de modo
    que las conjuntas y las condicionales inversas sigan visibles?
  \end{alertblock}

  \vspace{0.6em}
  Al final de esta lección, usted puede:
  \begin{itemize}
    \item dibujar un árbol de caminos terminales mutuamente excluyentes;
    \item etiquetar ramas con frecuencias o probabilidades;
    \item multiplicar a lo largo de un camino y sumar caminos disjuntos;
    \item explicar por qué $P(B\mid A)$ puede ser una rama mientras $P(A\mid B)$
      no lo es.
  \end{itemize}

  \vfill
  \scriptsize Todos los lotes y marcas de inspección usados hoy son sintéticos. No se usan datos reales de empresa.
\end{frame}

\begin{frame}{Un espacio muestral de dos etapas}
  \small
  \begin{center}
    \begin{tabular}{lrrr}
      \toprule
      \textbf{Estado del lote} & \textbf{Marcado} & \textbf{No marcado} & \textbf{Total} \\
      \midrule
      Defectuoso & 8 & 2 & 10 \\
      Limpio & 9 & 81 & 90 \\
      \midrule
      Total & 17 & 83 & 100 \\
      \bottomrule
    \end{tabular}
  \end{center}

  \vspace{0.5em}
  \begin{exampleblock}{Registro sintético de inspección}
    La etapa 1 es el estado del lote. La etapa 2 es el resultado de la
    inspección. Las cuatro celdas terminales son caminos mutuamente
    excluyentes que agotan el espacio muestral de 100 lotes.
  \end{exampleblock}
\end{frame}

\begin{frame}[fragile]{Paso 1: Árbol de frecuencias}
  \begin{columns}[T]
    \begin{column}{0.42\textwidth}
      \small
      \begin{lstlisting}
joints = {
  "def_flag": 8,
  "def_miss": 2,
  "clean_flag": 9,
  "clean_clear": 81,
}
flagged = 8 + 9
      \end{lstlisting}

      \begin{block}{Salida verificada}
        $100$ lotes; $10$ defectuosos\\
        Conteos terminales $8$, $2$, $9$, $81$\\
        Total marcado $=17$
      \end{block}
    \end{column}
    \begin{column}{0.58\textwidth}
      \centering
      \includegraphics[width=\textwidth]{../figuras/diagramas_arbol_01_arbol_frecuencias.png}
      \vspace{0.2em}
      \scriptsize Árbol de frecuencias del Paso 1
    \end{column}
  \end{columns}
\end{frame}

\begin{frame}{Multiplicar a lo largo de ramas; sumar caminos disjuntos}
  \small
  Ramas de primera etapa:
  \[
    P(D)=\frac{10}{100}=0.10,\qquad P(C)=\frac{90}{100}=0.90.
  \]
  Ramas de segunda etapa:
  \[
    P(F\mid D)=\frac{8}{10}=0.80,\qquad P(F\mid C)=\frac{9}{90}=0.10.
  \]
  Un producto de camino y el total marcado:
  \[
    P(D\cap F)=0.10\times 0.80=0.08,
  \]
  \[
    P(F)=0.08+0.09=0.17.
  \]
\end{frame}

\begin{frame}[fragile]{Paso 2: Probabilidades de camino}
  \begin{columns}[T]
    \begin{column}{0.40\textwidth}
      \small
      \begin{lstlisting}
p_def_flag = p_d * p_f_given_d
p_flag = (
    p_def_flag
    + p_clean_flag
)
      \end{lstlisting}

      \begin{block}{Salida verificada}
        $P(D\cap F)=0.080000$\\
        $P(C\cap F)=0.090000$\\
        $P(F)=0.170000$\\
        Suma de cuatro caminos $=1.000000$
      \end{block}
    \end{column}
    \begin{column}{0.60\textwidth}
      \centering
      \includegraphics[width=\textwidth]{../figuras/diagramas_arbol_02_probabilidades_caminos.png}
      \vspace{0.2em}
      \scriptsize Probabilidades de caminos terminales del Paso 2
    \end{column}
  \end{columns}
\end{frame}

\begin{frame}{Paso 3: La condicional inversa no es una rama}
  \begin{columns}[T]
    \begin{column}{0.42\textwidth}
      \small
      $P(F\mid D)=0.800000$ es una rama hacia abajo.

      \vspace{0.5em}
      \[
        P(D\mid F)=\frac{8}{17}=0.470588.
      \]

      \vspace{0.4em}
      Simulación semilla 42 de $100{,}000$ lotes: $0.469551$.

      \vspace{0.5em}
      \textbf{Límite:} un árbol en orden temporal no muestra $P(D\mid F)$
      como una sola rama. La Lección 11 nombra esa inversión como teorema de
      Bayes.
    \end{column}
    \begin{column}{0.58\textwidth}
      \centering
      \includegraphics[width=\textwidth]{../figuras/diagramas_arbol_03_condicional_inversa.png}
      \vspace{0.2em}
      \scriptsize Etiqueta de rama versus condicional inversa del Paso 3
    \end{column}
  \end{columns}
\end{frame}

\begin{frame}{Verificación de conceptos}
  \begin{enumerate}
    \item ¿Por qué los cuatro conteos terminales deben sumar 100 en este árbol?
    \item Calcule $P(C\cap F)$ a partir de las etiquetas de rama $0.90$ y
      $0.10$.
    \item ¿Cuál probabilidad es una sola rama de segunda etapa: $P(F\mid D)$ o
      $P(D\mid F)$?
    \item ¿Por qué $P(D\mid F)=8/17$ no es igual a $P(F\mid D)=0.80$?
  \end{enumerate}

  \vspace{0.6em}
  \begin{block}{Conexión con la Lección 11}
    La sesión puede continuar directamente con el teorema de Bayes,
    probabilidades priori y posterior, y ponderación por tasa base.
  \end{block}

  \vfill
  \scriptsize Fuente: OpenStax, \textit{Introductory Business Statistics 2e},
  Capítulo 3, Sección 3.4. CC BY-NC-SA.
\end{frame}

\appendix



\end{document}
